\chapter{Abordagem da Engenharia de Requisitos}

\section{Tradicional - RUP}
  Metodologia Tradicional $-$ RUP
            A metodologia tradicional nasceu com o crescimento da computação, onde o nascimento de sistemas empresariais e corporativos necessitou ferramentas para gerenciar grandes projetos. Uma das metodologias tradicionais mais utilizada é o RUP.
            RUP é um framework de processo de engenharia de software que fornece um conjunto de práticas testadas na indústria para desenvolvimento de software e gerência de projetos (Shuja, 2007). Como participante da categoria de metodologia tradicional, o RUP é caracterizado em fases e/ou etapas, e tem um foco alto na documentação do processo.
            Ele é uma metodologia iterativa, na qual basicamente é constituída em quatro fases e diversas iterações entre elas, na qual tem como finalidade tratar questões como planejamento, levantamento e análise de requisitos e entre outros. São elas:
          	\begin{description}
          	\item[$\bullet$] Iniciação: Tarefas de comunicação com o cliente e o planejamento;
          	\item[$\bullet$] Elaboração: Analisar de forma mais detalhada o domínio do problema;
          	\item[$\bullet$] Construção: Construção e desenvolvimento do software;
          	\item[$\bullet$] Transição: Entrega do software e testes.
          	\end{description}
 
No RUP também tem as disciplinas, que são basicamente atividades distribuídas ao longo das quatro fases, onde descrevem o que deve ser feito em cada fase com relação à documentação, atividades e outros.
 \begin{figure}[!htpb]
	\centering
	\includegraphics[scale=2.0]{figuras/abordagem/Fases_RUP}
	\caption{RUP}
\end{figure}
As principais características do RUP são:
\begin{description}
		 \item[$\bullet$] Casos de uso: Tem como finalidade capturar os requisitos funcionais do sistema;
         \item[$\bullet$] Centrado na arquitetura: Base do projeto deve ser sólida e estável, mas podendo ser flexível para mudanças e incrementos, ou seja, o sistema deverá ser o mais modularizado possível;
         \item[$\bullet$] Focado no risco: Controla os riscos mais críticos logo no início da fase de Elaboração;
         \item[$\bullet$] Baseado em componentes: Artefatos são construídos a partir da interconexão de objetos projetados através da linguagem UML.
\end{description}

Como o SAFe é baseado em processos ágeis, o framework também compartilha os princípios ágeis. Como por exemplo, o uso do kanban para os diferentes níveis, e até mesmo a utilização do scrum no nível mais baixo.

Além de compartilhar essas características com os processos ágeis, o SAFe possui 9 príncipos, são eles:
\begin{enumerate}
\item Ter uma visão econômica
\item Aplicar um pensamento sistêmico
\item Assumir variabilidade; preservar opções
\item Construir incrementalmente de forma rapida, com ciclos integrados de aprendizagem
\item Se baseie em marcos com o objetivo de avaliação de sistemas funcionais
\item Visualize e limite trabalhos em andamento, reduzir a quantidade de trabalhos e gerenciar grandes filas de espera.
\item Utilizar uma cadencia e sincronizar planejamento entre domínios
\item Habilitar a motivação intrínseca de conhecimento dos trabalhadores
\item Descentralizar a tomada de decisão
\end{enumerate}


O SAFe é divido em 3 níveis: \textbf{portfólio, programa e time.}


\subsection{Portfólio}
De acordo com o site do framework, o nível de portfólio é voltado para as pessoas e processos que são necessárias para a construção de sistemas e soluções que a empresa necessita para atingir uma meta prevista. 

\subsection{Programa}
O nível de programa é onde os times de desenvolvimento e outros recursos são aplicados em alguma solução em desenvolvimento corrente.

\subsection{Time}
Por fim, o nível de time mesmo que definido como níveis diferentes, é uma parte do nível de programa. Cada time é responsável por definir, construir e testar as histórias de seu backlog de time dentro de sprints, e cada um é organizado de acordo com as capacidades e talentos de cada indivíduo.

\subsection{Big picture}
\begin{figure}[!htpb]
	\centering
	\includegraphics[scale=0.7]{figuras/abordagem/SAFe_Big_Picture_4}
	\caption{Big picture do SAFe}
\end{figure}

\section{Tradicional}

\section{Híbrido}

\section{Justificativa}
